\documentclass[../../main.tex]{subfiles}
\graphicspath{{\subfix{../../image/}}}

\begin{document}

\section{强大数定律}\label{section:ch05-强大数定律}

\subsection*{一、以概率 $1$ 收敛}

以前，我们曾顺便提起过以概率 $1$ 收敛及强大数定律，本节将对它们进行深入讨论。

要彻底搞清以概率 $1$ 收敛这个概念，必须对事件（点集）序列的运算有进一步了解，我们就从讨论这个问题开始。

设 $A_1,A_2,\cdots,A_n,\cdots$ 是一列事件，则 $\bigcup_{n=k}^{\infty}A_n$ 表示事件序列 $A_k,A_{k+1},\cdots$ 中至少发生一个，而 $\bigcap_{n=k}^{\infty}A_n$ 则表示 $A_k,A_{k+1},\cdots$ 同时发生。

记
\begin{align}
\varlimsup_{n\to\infty}A_n=\bigcap_{k=1}^{\infty}\bigcup_{n=k}^{\infty}A_n,
\end{align}
\begin{align}
\varliminf_{n\to\infty}A_n=\bigcup_{k=1}^{\infty}\bigcap_{n=k}^{\infty}A_n.
\end{align}
称 $\varlimsup_{n\to\infty}A_n$ 为事件序列 $\{A_n\}$ 的\textbf{上限事件}，它表示 $A_n$ 发生无穷多次，因为 $\omega\in\bigcap_{k=1}^{\infty}\bigcup_{n=k}^{\infty}A_n$ 当且仅当 $\omega$ 属于无穷多个 $A_n$；类似地称 $\varliminf_{n\to\infty}A_n$ 为事件序列 $\{A_n\}$ 的\textbf{下限事件}，它表示 $A_n$ 至多只有有限个不发生，因为 $\omega\in\bigcup_{k=1}^{\infty}\bigcap_{n=k}^{\infty}A_n$，当且仅当存在一个 $N$，使 $\omega\in\bigcap_{n=N}^{\infty}A_n$，因此若 $\omega$ 发生，则 $A_N,A_{N+1},\cdots$ 同时发生，这时至多只有前面 $N-1$ 个事件 $A_1,A_2,\cdots,A_{N-1}$ 可能不发生（也可能有些发生）。

显然
\begin{align}
\varlimsup_{n\to\infty}A_n\supset\varliminf_{n\to\infty}A_n.
\end{align}
特别当 $\varlimsup_{n\to\infty}A_n=\varliminf_{n\to\infty}A_n$ 时，记 $\lim_{n\to\infty}A_n\equiv\varlimsup_{n\to\infty}A_n=\varliminf_{n\to\infty}A_n$，并称它为事件序列 $\{A_n\}$ 的\textbf{极限事件}。

利用 De Morgan 定理，有
\begin{align*}
\overline{\bigcap_{k=1}^{\infty}\bigcup_{n=k}^{\infty}A_n}=\bigcup_{k=1}^{\infty}\bigcap_{n=k}^{\infty}\bar{A}_n,
\end{align*}
\begin{align*}
\overline{\bigcup_{k=1}^{\infty}\bigcap_{n=k}^{\infty}A_n}=\bigcap_{k=1}^{\infty}\bigcup_{n=k}^{\infty}\bar{A}_n.
\end{align*}
因此
\begin{align}
\varliminf_{n\to\infty}\bar{A}_n=\overline{\varlimsup_{n\to\infty}A_n},
\end{align}
\begin{align}
\varlimsup_{n\to\infty}\bar{A}_n=\overline{\varliminf_{n\to\infty}A_n}.
\end{align}

下面 Borel--Cantelli 引理在概率论中有众多的应用。

\begin{lemma}[Borel--Cantelli 引理]\label{lemma:ch05-BorelCantelli}
\begin{enumerate}
\item 若随机事件序列 $\{A_n\}$ 满足
\begin{align}
\sum_{n=1}^{\infty}P(A_n)<\infty,
\end{align}
则
\begin{align}
P\{\varlimsup_{n\to\infty}A_n\}=0,\qquad P\{\varliminf_{n\to\infty}\bar{A}_n\}=1.
\end{align}
\item 若 $\{A_n\}$ 是相互独立的随机事件序列，则
\begin{align}
\sum_{n=1}^{\infty}P(A_n)=\infty
\end{align}
成立的充要条件为
\begin{align}
P\{\varlimsup_{n\to\infty}A_n\}=1\quad\text{或}\quad P\{\varliminf_{n\to\infty}\bar{A}_n\}=0.
\end{align}
\end{enumerate}
\end{lemma}

\begin{proof}
（i）由于
\begin{align*}
P\{\varlimsup_{n\to\infty}A_n\}=P\!\left\{\bigcap_{k=1}^{\infty}\bigcup_{n=k}^{\infty}A_n\right\}
\leqslant P\!\left\{\bigcup_{n=k}^{\infty}A_n\right\}\leqslant\sum_{n=k}^{\infty}P\{A_n\}\to0\qquad(k\to\infty).
\end{align*}
由上限下限事件的对偶关系
\begin{align*}
P\{\varliminf_{n\to\infty}\bar{A}_n\}=1.
\end{align*}

（ii）先证必要性。注意到 $\{A_n\}$ 的独立性，有
\begin{align*}
P\{\varliminf_{n\to\infty}\bar{A}_n\}=P\!\left\{\bigcup_{k=1}^{\infty}\bigcap_{n=k}^{\infty}\bar{A}_n\right\}
\leqslant\sum_{k=1}^{\infty}P\!\left\{\bigcap_{n=k}^{\infty}\bar{A}_n\right\}
=\sum_{k=1}^{\infty}\prod_{n=k}^{\infty}P(\bar{A}_n)
=\sum_{k=1}^{\infty}\prod_{n=k}^{\infty}[1-P(A_n)].
\end{align*}
由于
\begin{align*}
0\leqslant1-P(A_n)\leqslant\exp\{-P(A_n)\},
\end{align*}
则从
\begin{align*}
\sum_{n=1}^{\infty}P(A_n)=\infty
\end{align*}
可得
\begin{align*}
\prod_{n=k}^{\infty}[1-P(A_n)]\leqslant\lim_{N\to\infty}\exp\!\left\{-\sum_{n=k}^{N}P(A_n)\right\}=0.
\end{align*}
所以
\begin{align*}
P\{\varliminf_{n\to\infty}\bar{A}_n\}=0.
\end{align*}

再证充分性。若 $P\{\varlimsup_{n\to\infty}A_n\}=1$。假定 $\sum_{n=1}^{\infty}P(A_n)<\infty$，则由（i）得到 $P\{\varlimsup_{n\to\infty}A_n\}=0$，产生矛盾。因 $P(A_n)\geqslant0$，故只可能是
\begin{align*}
\sum_{n=1}^{\infty}P(A_n)=\infty,
\end{align*}
引理证毕。
\end{proof}

现在讨论随机变量序列的以概率 $1$ 收敛性。

若 $\xi_n(\omega)$（$n=1,2,\cdots$），$\xi(\omega)$ 是随机变量，则
\begin{align}
\{\omega:\lim_{n\to\infty}\xi_n(\omega)=\xi(\omega)\}
=\left\{\omega:\bigcap_{m=1}^{\infty}\bigcup_{k=1}^{\infty}\bigcap_{n=k}^{\infty}\left(|\xi_n(\omega)-\xi(\omega)|<\frac1m\right)\right\}.
\end{align}
这个式子可以这样理解：因为 $\omega\in\{\lim_{n\to\infty}\xi_n(\omega)=\xi(\omega)\}$ 的充要条件是：对任一正整数 $m$，存在一个正整数 $N$，使当 $n>N$ 时均有 $|\xi_n(\omega)-\xi(\omega)|<\dfrac1m$；即对任一正整数 $m$，$\omega$ 属于 $\left(|\xi_n(\omega)-\xi(\omega)|<\dfrac1m\right)$ 的下限事件，这正是上式的右边。从这个表达式中还可以看出，$\{\lim_{n\to\infty}\xi_n(\omega)=\xi(\omega)\}$ 是事件，因此
\begin{align}
P\{\lim_{n\to\infty}\xi_n(\omega)=\xi(\omega)\}=1
\end{align}
有明确的意义，这时称 $\{\xi_n(\omega)\}$ \textbf{以概率 $1$ 收敛}于 $\xi(\omega)$，记为 $\xi_n(\omega)\xrightarrow{\mathrm{a.s.}}\xi(\omega)$。

因此下面两个式子都表达了 $\{\xi_n(\omega)\}$ 以概率 $1$ 收敛于 $\xi(\omega)$。
\begin{align}
P\!\left\{\bigcap_{m=1}^{\infty}\bigcup_{k=1}^{\infty}\bigcap_{n=k}^{\infty}\left(|\xi_n(\omega)-\xi(\omega)|<\frac1m\right)\right\}=1,
\end{align}
\begin{align}
P\!\left\{\bigcup_{m=1}^{\infty}\bigcap_{k=1}^{\infty}\bigcup_{n=k}^{\infty}\left(|\xi_n(\omega)-\xi(\omega)|\geqslant\frac1m\right)\right\}=0.
\end{align}

进一步，我们要说明，$\{\xi_n(\omega)\}$ 以概率 $1$ 收敛于 $\xi(\omega)$ 的定义也可以表达为：对任意的 $\varepsilon>0$，成立
\begin{align}
P\!\left\{\bigcap_{k=1}^{\infty}\bigcup_{n=k}^{\infty}\left(|\xi_n(\omega)-\xi(\omega)|\geqslant\varepsilon\right)\right\}=0.
\end{align}
若以 $A_n$ 记 $(|\xi_n(\omega)-\xi(\omega)|\geqslant\varepsilon)$，上式表示 $\varlimsup_{n\to\infty}A_n$ 的概率为 $0$，这与我们对以概率 $1$ 收敛的理解一致。不过，对于这个结论，还是给它一个严格的证明。

事实上，由于对 $\varepsilon>0$，总有
\begin{align*}
\left\{\bigcap_{k=1}^{\infty}\bigcup_{n=k}^{\infty}\left(|\xi_n(\omega)-\xi(\omega)|\geqslant\varepsilon\right)\right\}
\subset\left\{\bigcup_{m=1}^{\infty}\bigcap_{k=1}^{\infty}\bigcup_{n=k}^{\infty}\left(|\xi_n(\omega)-\xi(\omega)|\geqslant\frac1m\right)\right\},
\end{align*}
因此由前述等式可以推得极限为 $0$ 的结论。反之，利用
\begin{align*}
P\!\left\{\bigcup_{m=1}^{\infty}\bigcap_{k=1}^{\infty}\bigcup_{n=k}^{\infty}\left(|\xi_n(\omega)-\xi(\omega)|\geqslant\frac1m\right)\right\}
\leqslant\sum_{m=1}^{\infty}P\!\left\{\bigcap_{k=1}^{\infty}\bigcup_{n=k}^{\infty}\left(|\xi_n(\omega)-\xi(\omega)|\geqslant\frac1m\right)\right\},
\end{align*}
可由极限为 $0$ 的结论推出前式，这就说明了两种表达法的等价性。

利用概率的连续性可知，该式等价于
\begin{align}
\lim_{k\to\infty}P\!\left\{\bigcup_{n=k}^{\infty}\left(|\xi_n(\omega)-\xi(\omega)|\geqslant\varepsilon\right)\right\}=0.
\end{align}
根据 De Morgan 定理又知该式等价于
\begin{align}
\lim_{k\to\infty}P\!\left\{\bigcap_{n=k}^{\infty}\left(|\xi_n(\omega)-\xi(\omega)|<\varepsilon\right)\right\}=1.
\end{align}
由于
\begin{align*}
|\xi_k(\omega)-\xi(\omega)|\geqslant\varepsilon\subset\left\{\bigcup_{n=k}^{\infty}\left(|\xi_n(\omega)-\xi(\omega)|\geqslant\varepsilon\right)\right\},
\end{align*}
因此若该式成立，则
\begin{align*}
\lim_{k\to\infty}P\{|\xi_k(\omega)-\xi(\omega)|\geqslant\varepsilon\}=0.
\end{align*}
这样一来，我们已证得

\begin{theorem}\label{theorem:ch05-a.s.推依概率}
$\xi_n(\omega)\xrightarrow{\mathrm{a.s.}}\xi(\omega)\Rightarrow\xi_n(\omega)\xrightarrow{P}\xi(\omega)$。
\end{theorem}

下例说明一般不能由依概率收敛推得以概率 $1$ 收敛，所以以概率 $1$ 收敛是比依概率收敛更强的一种收敛性。

\begin{example}\label{example:ch05-依概率不a.s.反例}
取 $\Omega=(0,1]$，$\mathscr{F}$ 为 $(0,1]$ 中 Borel 点集全体所构成的 $\sigma$ 域，$P$ 为 Lebesgue 测度，令
\begin{align}
\eta_{ki}(\omega)=
\begin{cases}
1,&\omega\in\left(\dfrac{i-1}{k},\dfrac{i}{k}\right],\quad i=1,2,\cdots,k,\\
0,&\omega\bar{\in}\left(\dfrac{i-1}{k},\dfrac{i}{k}\right],\quad k=1,2,\cdots.
\end{cases}
\end{align}
定义
\begin{align*}
\xi_1(\omega)=\eta_{11}(\omega),\xi_2(\omega)=\eta_{21}(\omega),\xi_3(\omega)=\eta_{22}(\omega),\\
\xi_4(\omega)=\eta_{31}(\omega),\xi_5(\omega)=\eta_{32}(\omega),\cdots
\end{align*}
一般 $\xi_n(\omega)=\eta_{ki}(\omega)$，其中 $n=i+\dfrac{k(k-1)}2$。这样定义的 $\{\xi_n(\omega)\}$ 是一列随机变量。但对于任何一个 $\omega\in(0,1]$，$\xi_n(\omega)$ 必有无限个 $k,i$ 使其取值 $0$，也有无限个 $k,i$ 使其取值 $1$，因此 $\{\xi_n(\omega)\}$ 不是以概率 $1$ 收敛于 $0$。但是另一方面，对任意的 $\varepsilon>0$，
\begin{align*}
P\{|\eta_{ki}(\omega)|\geqslant\varepsilon\}\leqslant\frac1k.
\end{align*}
当 $n\to\infty$ 时，由 $n=\dfrac{k(k-1)}2+i\leqslant\dfrac{k(k-1)}2+k$，知道 $k\to\infty$，因此
\begin{align*}
\lim_{n\to\infty}P\{|\xi_n(\omega)|\geqslant\varepsilon\}
=\lim_{n\to\infty}P\{|\eta_{ki}(\omega)|\geqslant\varepsilon\}=0.
\end{align*}
所以 $\{\xi_n(\omega)\}$ 依概率收敛于 $0$。

不难验证，$\{\xi_n(\omega)\}$ 是 $r$ 阶收敛于 $0$ 的，因此本例也提供了 $r$ 阶收敛推不出以概率 $1$ 收敛之例。
\end{example}

我们以前讨论的大数定律只要求依概率收敛，若把收敛性要求提高为以概率 $1$ 收敛，则得到的大数定律称为\textbf{强大数定律}（strong law of large numbers）。由\cref{theorem:ch05-a.s.推依概率}可知，若强大数定律成立，则通常的大数定律也一定成立，反之不然。有时为区别起见，把依概率收敛意义下的大数定律称为\textbf{弱大数定律}（weak law of large numbers）。

第一个强大数定律是由 Borel 在 1909 年对 Bernoulli 试验场合建立的。

\subsection*{二、Borel 强大数定律}

\begin{theorem}[Borel]\label{theorem:ch05-Borel强大数定律}
设 $\mu_n$ 是事件 $A$ 在 $n$ 次独立试验中的出现次数，在每次试验中事件 $A$ 出现的概率均为 $p$，那么当 $n\to\infty$ 时，
\begin{align}
P\!\left\{\frac{\mu_n}{n}\to p\right\}=1.
\end{align}
\end{theorem}

\begin{proof}
为使上式成立，由以概率 $1$ 收敛的等价条件知，只须对任意的 $\varepsilon>0$，成立
\begin{align}
P\!\left\{\bigcap_{k=1}^{\infty}\bigcup_{n=k}^{\infty}\left(\left|\frac{\mu_n}{n}-p\right|\geqslant\varepsilon\right)\right\}=0.
\end{align}
若记 $A_n=\left\{\left|\dfrac{\mu_n}{n}-p\right|\geqslant\varepsilon\right\}$，则上式可写成 $P\{\varlimsup_{n\to\infty}A_n\}=0$。根据 Borel--Cantelli 引理，为证明该式只要能证明级数
\begin{align}
\sum_{n=1}^{\infty}P\!\left\{\left|\frac{\mu_n}{n}-p\right|\geqslant\varepsilon\right\}
\end{align}
对任何 $\varepsilon>0$ 都收敛就可以了。

假如像证明 Bernoulli 大数定律那样用 Chebyshev 不等式进行估计，只能得到
\begin{align}
P\!\left\{\left|\frac{\mu_n}{n}-p\right|\geqslant\varepsilon\right\}\leqslant\frac{1}{4n\varepsilon^{2}}.
\end{align}
这对证明弱大数定律足够了，但为了保证级数收敛还不行，这时必须寻找更好的估计式。在这种特殊场合，Markov 不等式就够用了。由于
\begin{align}
P\!\left\{\left|\frac{\mu_n}{n}-p\right|\geqslant\varepsilon\right\}\leqslant\frac{1}{\varepsilon^{4}}\boldsymbol{E}\left|\frac{\mu_n}{n}-p\right|^{4}.
\end{align}
问题是要计算 $\dfrac{\mu_n}{n}$ 的四阶中心矩。还是像过去一样，我们把 $\mu_n$ 表示成独立 Bernoulli $0$--$1$ 变量 $\xi_1,\xi_2,\cdots,\xi_n$ 之和，这样
\begin{align*}
\frac{\mu_n}{n}-p=\frac1n\sum_{i=1}^{n}(\xi_i-p).
\end{align*}
所以
\begin{align*}
\boldsymbol{E}\!\left(\frac{\mu_n}{n}-p\right)^{4}
=\frac{1}{n^{4}}\sum_{i=1}^{n}\sum_{j=1}^{n}\sum_{k=1}^{n}\sum_{l=1}^{n}\boldsymbol{E}(\xi_i-p)(\xi_j-p)(\xi_k-p)(\xi_l-p).
\end{align*}
注意到各 $\xi_i$ 的独立性及 $\boldsymbol{E}(\xi_i-p)=0$，因此上面的和式中只有 $\boldsymbol{E}(\xi_i-p)^{4}$ 及 $\boldsymbol{E}(\xi_i-p)^{2}(\xi_j-p)^{2}$ 的项才不等于 $0$，显然
\begin{align}
\boldsymbol{E}(\xi_i-p)^{4}=pq(p^{3}+q^{3}),
\end{align}
\begin{align}
\boldsymbol{E}(\xi_i-p)^{2}(\xi_j-p)^{2}=p^{2}q^{2}\qquad(i\neq j).
\end{align}
四阶矩形式的项有 $n$ 项，二阶矩形式的项有 $\binom{4}{2}\binom{n}{2}=3n(n-1)$ 项，因此
\begin{align}
\boldsymbol{E}\!\left(\frac{\mu_n}{n}-p\right)^{4}
=\frac{pq}{n^{4}}\bigl[n(p^{3}+q^{3})+3pq(n^{2}-n)\bigr]<\frac{1}{4n^{2}}.
\end{align}
于是
\begin{align}
P\!\left\{\left|\frac{\mu_n}{n}-p\right|\geqslant\varepsilon\right\}<\frac{1}{4\varepsilon^{4}n^{2}}.
\end{align}
这个估计式已经比 Chebyshev 估计进了一大步，它可以保证级数收敛，从而证明了定理。
\end{proof}

从本书第一节介绍随机事件频率稳定性时，我们就期待着这样一个结论，即当试验次数无限增加时，频率将趋于概率，Borel 强大数定律正给出了这个结果。从 Bernoulli 大数定律并不能引申出这个结论，它只断言一个不等式 $\left|\dfrac{\mu_n}{n}-p\right|<\varepsilon$ 成立的概率可以大于 $1-\eta$，不论 $\eta$ 是什么正数；但是事件
\begin{align*}
\left|\frac{\mu_{n+1}}{n+1}-p\right|\geqslant\varepsilon,\left|\frac{\mu_{n+2}}{n+2}-p\right|\geqslant\varepsilon,\cdots,\left|\frac{\mu_{2n}}{2n}-p\right|\geqslant\varepsilon,\cdots
\end{align*}
中至少有一个发生仍是可能的，因为它是可列个事件之并，而我们只知道每个事件的概率很小。但 Borel 强大数定律则断言 $\left|\dfrac{\mu_n}{n}-p\right|$ 以概率 $1$ 变得很小，而且保持很小。虽然从逻辑上讲，在投硬币时每次都出现正面是可能的，这时 $\dfrac{\mu_n}{n}=1$，因而 $\dfrac{\mu_n}{n}\to p$ 并不成立，但是强大数定律断言了这种事件发生的概率为 $0$。

\subsection*{三、Kolmogorov 强大数定律}

下面讨论更一般的强大数定律，先把其含义进一步明确如下：

设 $\{\xi_i\}$ 是独立随机变量序列，若
\begin{align}
P\!\left\{\lim_{n\to\infty}\frac1n\sum_{i=1}^{n}(\xi_i-\boldsymbol{E}\xi_i)=0\right\}=1
\end{align}
则称它满足强大数定律。

根据以概率 $1$ 收敛的等价条件，这等价于要求对任意 $\varepsilon>0$ 成立
\begin{align}
\lim_{m\to\infty}P\!\left\{\bigcup_{j=m}^{\infty}\left(\left|\frac1j\sum_{i=1}^{j}(\xi_i-\boldsymbol{E}\xi_i)\right|\geqslant\varepsilon\right)\right\}=0.
\end{align}
由于
\begin{align}
\bigcup_{j=m}^{\infty}\left(\left|\frac1j\sum_{i=1}^{j}(\xi_i-\boldsymbol{E}\xi_i)\right|\geqslant\varepsilon\right)
\subset\left\{\sup_{j\geqslant m}\left|\frac1j\sum_{i=1}^{j}(\xi_i-\boldsymbol{E}\xi_i)\right|\geqslant\varepsilon\right\},
\end{align}
因此需要对概率 $P\left\{\sup_{j\geqslant m}\left|\frac1j\sum_{i=1}^{j}(\xi_i-\boldsymbol{E}\xi_i)\right|\geqslant\varepsilon\right\}$ 进行估计，这相当于在独立和场合对 Chebyshev 不等式进行推广。这方面已经有不少成果，在这里我们介绍一个由 Hájek 及 Rényi 证明的不等式。

\begin{theorem}[Hájek--Rényi 不等式]\label{theorem:ch05-HajekRenyi}
若 $\{\xi_i\}$ 是独立随机变量序列，$\boldsymbol{D}\xi_i=\sigma_i^{2}<\infty$（$i=1,2,\cdots$），而 $\{C_n\}$ 是一列正的非增常数序列，则对任意正整数 $m,n$（$m<n$）及 $\varepsilon>0$，均有
\begin{align}
P\!\left\{\max_{m\leqslant j\leqslant n}C_j\left|\sum_{i=1}^{j}(\xi_i-\boldsymbol{E}\xi_i)\right|\geqslant\varepsilon\right\}
\leqslant\frac{1}{\varepsilon^{2}}\left(C_m^{2}\sum_{j=1}^{m}\sigma_j^{2}+\sum_{j=m+1}^{n}C_j^{2}\sigma_j^{2}\right).
\end{align}
\end{theorem}

\begin{proof}
记
\begin{align*}
S_k=\sum_{j=1}^{k}(\xi_j-\boldsymbol{E}\xi_j)
\end{align*}
及
\begin{align}
\eta=\sum_{k=m}^{n-1}S_k^{2}(C_k^{2}-C_{k+1}^{2})+C_n^{2}S_n^{2}.
\end{align}
因此
\begin{align*}
\eta&=\sum_{k=m}^{n}S_k^{2}C_k^{2}-\sum_{k=m}^{n-1}S_k^{2}C_{k+1}^{2}\\
&=S_m^{2}C_m^{2}+\sum_{k=m+1}^{n}(S_k^{2}-S_{k-1}^{2})C_k^{2}.
\end{align*}
利用 $\{\xi_i\}$ 的独立性
\begin{align*}
\boldsymbol{E}S_k^{2}=\sum_{j=1}^{k}\sigma_j^{2},
\end{align*}
所以有
\begin{align}
\boldsymbol{E}\eta=C_m^{2}\sum_{j=1}^{m}\sigma_j^{2}+\sum_{j=m+1}^{n}C_j^{2}\sigma_j^{2}.
\end{align}
对 $j=m,m+1,\cdots,n$，记
\begin{align}
E_j=\{C_k|S_k|<\varepsilon,m\leqslant k<j;C_j|S_j|\geqslant\varepsilon\}.
\end{align}
这样定义的 $E_j$（$j=m,m+1,\cdots,n$）是互不相容的，而且
\begin{align}
P\!\left\{\max_{m\leqslant j\leqslant n}C_j\left|\sum_{i=1}^{j}(\xi_i-\boldsymbol{E}\xi_i)\right|\geqslant\varepsilon\right\}
=P\!\left\{\max_{m\leqslant j\leqslant n}C_j|S_j|\geqslant\varepsilon\right\}
=\sum_{j=m}^{n}P(E_j).
\end{align}
令
\begin{align}
\chi_j(\omega)=
\begin{cases}
1,&\omega\in E_j,\\
0,&\omega\bar{\in}E_j,
\end{cases}\quad j=m,m+1,\cdots,n.
\end{align}
即 $\chi_j$ 是 $E_j$ 的示性函数，注意到 $E_j$ 的互不相容性及 $\sum_{j=m}^{n}E_j\subset\Omega$，因而
\begin{align}
\sum_{j=m}^{n}\chi_j(\omega)\leqslant1.
\end{align}
所以有
\begin{align}
\boldsymbol{E}\eta\geqslant\sum_{j=m}^{n}\boldsymbol{E}(\eta\chi_j).
\end{align}
对 $j<k\leqslant n$，
\begin{align*}
S_k=S_j+(\xi_{j+1}-\boldsymbol{E}\xi_{j+1})+\cdots+(\xi_k-\boldsymbol{E}\xi_k).
\end{align*}
因此
\begin{align}
\boldsymbol{E}(S_k^{2}\chi_j)&=\boldsymbol{E}(S_j^{2}\chi_j)\notag\\
&\quad+\boldsymbol{E}\left\{\bigl[(\xi_{j+1}-\boldsymbol{E}\xi_{j+1})+\cdots+(\xi_k-\boldsymbol{E}\xi_k)\bigr]^{2}\cdot\chi_j\right\}\notag\\
&\quad+2\boldsymbol{E}\left\{S_j\bigl[(\xi_{j+1}-\boldsymbol{E}\xi_{j+1})+\cdots+(\xi_k-\boldsymbol{E}\xi_k)\bigr]\cdot\chi_j\right\}\notag\\
&\geqslant\boldsymbol{E}(S_j^{2}\chi_j)+2\boldsymbol{E}\{S_j\chi_j[(\xi_{j+1}-\boldsymbol{E}\xi_{j+1})+\notag\\
&\qquad\qquad\qquad\cdots+(\xi_k-\boldsymbol{E}\xi_k)]\}.
\end{align}
由于 $S_j\chi_j$ 只与 $\xi_1,\cdots,\xi_j$ 有关，因此与 $[(\xi_{j+1}-\boldsymbol{E}\xi_{j+1})+\cdots+(\xi_k-\boldsymbol{E}\xi_k)]$ 独立，故
\begin{align}
\boldsymbol{E}\{S_j\chi_j[(\xi_{j+1}-\boldsymbol{E}\xi_{j+1})+\cdots+(\xi_k-\boldsymbol{E}\xi_k)]\}
=\boldsymbol{E}(S_j\chi_j)\cdot\boldsymbol{E}[(\xi_{j+1}-\boldsymbol{E}\xi_{j+1})+\cdots+(\xi_k-\boldsymbol{E}\xi_k)]=0.
\end{align}
而在 $E_j$ 上 $|S_j|\geqslant\dfrac{\varepsilon}{C_j}$，故有
\begin{align}
\boldsymbol{E}(S_j^{2}\chi_j)\geqslant\frac{\varepsilon^{2}}{C_j^{2}}\boldsymbol{E}\chi_j=\frac{\varepsilon^{2}}{C_j^{2}}P(E_j).
\end{align}
因此由前述几式得到
\begin{align}
\boldsymbol{E}(S_k^{2}\chi_j)\geqslant\frac{\varepsilon^{2}}{C_j^{2}}P(E_j),\quad j\leqslant k\leqslant n.
\end{align}
现在当 $m\leqslant j\leqslant n$ 时，由 $\eta$ 的定义知
\begin{align}
\boldsymbol{E}(\eta\chi_j)&=\sum_{k=m}^{n-1}\boldsymbol{E}(S_k^{2}\chi_j)(C_k^{2}-C_{k+1}^{2})+C_n^{2}\boldsymbol{E}(S_n^{2}\chi_j)\notag\\
&\geqslant\sum_{k=m}^{n-1}\boldsymbol{E}(S_k^{2}\chi_j)(C_k^{2}-C_{k+1}^{2})+C_n^{2}\boldsymbol{E}(S_n^{2}\chi_j)\notag\\
&\geqslant\frac{\varepsilon^{2}}{C_j^{2}}P(E_j)\left[\sum_{k=j}^{n-1}(C_k^{2}-C_{k+1}^{2})+C_n^{2}\right]\notag\\
&=\varepsilon^{2}P(E_j).
\end{align}
由 $\boldsymbol{E}\eta\geqslant\sum_{j=m}^{n}\boldsymbol{E}(\eta\chi_j)$ 及上式
\begin{align}
\boldsymbol{E}\eta\geqslant\varepsilon^{2}\sum_{j=m}^{n}P(E_j).
\end{align}
利用 $\boldsymbol{E}\eta$ 的表达式、概率等式及上式即得\cref{theorem:ch05-HajekRenyi}。
\end{proof}

在 Hájek--Rényi 不等式中，特别令 $m=1$，$C_j=1$，则得到著名的 Kolmogorov 不等式。

\begin{theorem}[Kolmogorov 不等式]\label{theorem:ch05-Kolmogorov不等式}
设 $\xi_1,\xi_2,\cdots,\xi_n$ 是独立随机变量，方差有限，则对任意 $\varepsilon>0$，成立
\begin{align}
P\!\left\{\max_{1\leqslant j\leqslant n}\left|\sum_{i=1}^{j}(\xi_i-\boldsymbol{E}\xi_i)\right|\geqslant\varepsilon\right\}
\leqslant\frac{1}{\varepsilon^{2}}\sum_{j=1}^{n}\boldsymbol{D}\xi_j.
\end{align}
\end{theorem}

Kolmogorov 不等式是概率论中最重要的不等式之一，有广泛的应用。在上式中，若令 $n=1$，则得到
\begin{align*}
P\{|\xi_1-\boldsymbol{E}\xi_1|\geqslant\varepsilon\}\leqslant\frac{\boldsymbol{D}\xi_1}{\varepsilon^{2}},
\end{align*}
这正是 Chebyshev 不等式。因此 Kolmogorov 不等式是 Chebyshev 不等式的推广，而 Hájek--Rényi 不等式又是 Kolmogorov 不等式的推广。

利用 Hájek--Rényi 不等式，能证明下面重要结果。

\begin{theorem}[Kolmogorov 强大数定律]\label{theorem:ch05-Kolmogorov强大数定律}
设 $\{\xi_i\},i=1,2,\cdots$ 是独立随机变量序列，且 $\sum_{n=1}^{\infty}\dfrac{\boldsymbol{D}\xi_n}{n^{2}}<\infty$，则成立
\begin{align}
P\!\left\{\lim_{n\to\infty}\frac1n\sum_{i=1}^{n}(\xi_i-\boldsymbol{E}\xi_i)=0\right\}=1.
\end{align}
\end{theorem}

\begin{proof}
在 Hájek--Rényi 不等式中，令 $C_j=\dfrac1j$，可以得到
\begin{align*}
P\!\left\{\max_{m\leqslant j\leqslant n}\left|\frac1j\sum_{i=1}^{j}(\xi_i-\boldsymbol{E}\xi_i)\right|\geqslant\varepsilon\right\}
\leqslant\frac{1}{\varepsilon^{2}}\left(\frac{1}{m^{2}}\sum_{j=1}^{m}\boldsymbol{D}\xi_j+\sum_{j=m+1}^{n}\frac{\boldsymbol{D}\xi_j}{j^{2}}\right).
\end{align*}
由概率的连续性
\begin{align}
P\!\left\{\sup_{j\geqslant m}\left|\frac1j\sum_{i=1}^{j}(\xi_i-\boldsymbol{E}\xi_i)\right|\geqslant\varepsilon\right\}
&=\lim_{n\to\infty}P\!\left\{\max_{m\leqslant j\leqslant n}\left|\frac1j\sum_{i=1}^{j}(\xi_i-\boldsymbol{E}\xi_i)\right|\geqslant\varepsilon\right\}\notag\\
&\leqslant\frac{1}{\varepsilon^{2}}\left(\frac{1}{m^{2}}\sum_{j=1}^{m}\boldsymbol{D}\xi_j+\sum_{j=m+1}^{\infty}\frac{\boldsymbol{D}\xi_j}{j^{2}}\right).
\end{align}
因为 $\sum_{j=1}^{\infty}\dfrac{\boldsymbol{D}\xi_j}{j^{2}}<\infty$，故由上式得到
\begin{align*}
\lim_{m\to\infty}P\!\left\{\sup_{j\geqslant m}\left|\frac1j\sum_{i=1}^{j}(\xi_i-\boldsymbol{E}\xi_i)\right|\geqslant\varepsilon\right\}=0.
\end{align*}
再由上限事件包含关系知强大数定律成立。定理证毕。

显然，由 Kolmogorov 强大数定律很容易推出 Borel 强大数定律。
\end{proof}

\subsection*{四、独立同分布场合的强大数定律}

在这种特殊的场合，可以找到强大数定律成立的充要条件，这个结果也属于 Kolmogorov。

\begin{theorem}[Kolmogorov]\label{theorem:ch05-独立同分布强大数定律}
设 $\xi_1,\xi_2,\cdots$ 是相互独立相同分布的随机变量序列，则
\begin{align}
\frac1n(\xi_1+\xi_2+\cdots+\xi_n)\xrightarrow{\mathrm{a.s.}}a
\end{align}
成立的充要条件是 $\boldsymbol{E}\xi_i$ 存在且等于 $a$。
\end{theorem}

\begin{proof}
若 $\xi$ 的分布函数为 $F(x)$，我们来证明不等式：
\begin{align}
\sum_{n=1}^{\infty}P\{|\xi|\geqslant n\}\leqslant\boldsymbol{E}|\xi|\leqslant1+\sum_{n=1}^{\infty}P\{|\xi|\geqslant n\}.
\end{align}
事实上，
\begin{align*}
\boldsymbol{E}|\xi|=\int_{-\infty}^{\infty}|x|\,\mathrm{d}F(x)=\sum_{k=0}^{\infty}\int_{k\leqslant|x|<k+1}|x|\,\mathrm{d}F(x).
\end{align*}
因此
\begin{align*}
\sum_{k=0}^{\infty}kP\{k\leqslant|\xi|<k+1\}\leqslant\boldsymbol{E}|\xi|
\leqslant\sum_{k=0}^{\infty}(k+1)P\{k\leqslant|\xi|<k+1\}.
\end{align*}
现在有
\begin{align*}
\sum_{k=0}^{\infty}kP\{k\leqslant|\xi|<k+1\}
=\sum_{n=1}^{\infty}\sum_{k=n}^{\infty}P\{k\leqslant|\xi|<k+1\}
=\sum_{n=1}^{\infty}P\{|\xi|\geqslant n\},
\end{align*}
及
\begin{align*}
\sum_{k=0}^{\infty}(k+1)P\{k\leqslant|\xi|<k+1\}
=\sum_{k=0}^{\infty}kP\{k\leqslant|\xi|<k+1\}+1
=\sum_{n=1}^{\infty}P\{|\xi|\geqslant n\}+1.
\end{align*}
这就证得了上述不等式。这个不等式说明 $\boldsymbol{E}|\xi|<\infty$ 的充要条件为
\begin{align*}
\sum_{n=1}^{\infty}P\{|\xi|\geqslant n\}<\infty.
\end{align*}

记 $S_n=\xi_1+\xi_2+\cdots+\xi_n$，若 $\dfrac{S_n}{n}\xrightarrow{\mathrm{a.s.}}\mu$，这里 $\mu$ 是有限数，则
\begin{align}
\frac{\xi_n}{n}=\frac{S_n}{n}-\frac{n-1}{n}\cdot\frac{S_{n-1}}{n-1}\xrightarrow{\mathrm{a.s.}}0.
\end{align}
这样一来，事件 $\{|\xi_n|\geqslant n\}$ 发生无穷多次的概率为 $0$，因此注意到 $\xi_i$ 的独立性，并利用 Borel--Cantelli 引理（ii），可知
\begin{align*}
\sum_{n=1}^{\infty}P\{|\xi_n|\geqslant n\}<\infty.
\end{align*}
再由上述不等式即知 $\boldsymbol{E}|\xi_i|<\infty$，这时显然有 $a=\boldsymbol{E}\xi_i$。这样，我们已证得必要性。

下证充分性。用``截尾法''，令
\begin{align}
\xi_n^{*}=
\begin{cases}
\xi_n,&|\xi_n|<n,\\
0,&|\xi_n|\geqslant n.
\end{cases}
\end{align}
先验证 $\{\xi_n^{*}\}$ 满足 Kolmogorov 强大数定律条件。以 $F(x)$ 记 $\xi_i$ 的分布函数，则
\begin{align*}
\boldsymbol{D}\xi_n^{*}\leqslant\boldsymbol{E}\xi_n^{*2}=\int_{-n}^{n}x^{2}\,\mathrm{d}F(x)\leqslant\sum_{k=1}^{n}k^{2}P\{k-1\leqslant|\xi_n|<k\},
\end{align*}
\begin{align*}
\sum_{n=1}^{\infty}\frac{\boldsymbol{D}\xi_n^{*}}{n^{2}}
\leqslant\sum_{n=1}^{\infty}\sum_{k=1}^{n}\frac{k^{2}}{n^{2}}P\{k-1\leqslant|\xi_n|<k\}
=\sum_{k=1}^{\infty}\sum_{n=k}^{\infty}\frac{k^{2}}{n^{2}}P\{k-1\leqslant|\xi_n|<k\}.
\end{align*}
由于
\begin{align*}
\sum_{n=k}^{\infty}\frac{1}{n^{2}}<\frac{1}{k^{2}}+\sum_{n=k+1}^{\infty}\frac{1}{n(n-1)}=\frac{1}{k^{2}}+\frac1k\leqslant\frac{2}{k},
\end{align*}
故
\begin{align*}
\sum_{n=1}^{\infty}\frac{\boldsymbol{D}\xi_n^{*}}{n^{2}}<2\cdot\sum_{k=1}^{\infty}kP\{k-1\leqslant|\xi_n|<k\}<\infty.
\end{align*}
因此
\begin{align}
P\!\left\{\lim_{n\to\infty}\frac1n\sum_{i=1}^{n}(\xi_i^{*}-\boldsymbol{E}\xi_i^{*})=0\right\}=1.
\end{align}
因为
\begin{align*}
\boldsymbol{E}\xi_n^{*}=\int_{-n}^{n}x\,\mathrm{d}F(x),
\end{align*}
显然 $\lim_{n\to\infty}\boldsymbol{E}\xi_n^{*}=\boldsymbol{E}\xi_1=a$，因此 $\lim_{n\to\infty}\frac1n\sum_{i=1}^{n}\boldsymbol{E}\xi_i^{*}=\boldsymbol{E}\xi_1=a$。由于
\begin{align}
\left|\frac1n\sum_{i=1}^{n}(\xi_i-a)\right|
\leqslant\left|\frac1n\sum_{i=1}^{n}(\xi_i-\xi_i^{*})\right|
+\left|\frac1n\sum_{i=1}^{n}(\xi_i^{*}-\boldsymbol{E}\xi_i^{*})\right|
+\left|\frac1n\sum_{i=1}^{n}(\boldsymbol{E}\xi_i^{*}-a)\right|.
\end{align}
为证 $\frac1n(\xi_1+\cdots+\xi_n)\xrightarrow{\mathrm{a.s.}}a$ 成立，只须再证 $\frac1n\sum_{i=1}^{n}(\xi_i-\xi_i^{*})\xrightarrow{\mathrm{a.s.}}0$。然而
\begin{align*}
\sum_{i=1}^{\infty}P\{\xi_i\neq\xi_i^{*}\}=\sum_{i=1}^{\infty}P\{|\xi_i|\geqslant i\}\leqslant\boldsymbol{E}|\xi_1|<\infty.
\end{align*}
由 Borel--Cantelli 引理知，以概率 $1$ 有
\begin{align*}
\xi_i(\omega)\neq\xi_i^{*}(\omega),\text{只对有限个 }i\text{ 成立},
\end{align*}
因此
\begin{align*}
P\!\left\{\lim_{n\to\infty}\frac1n\sum_{i=1}^{n}(\xi_i-\xi_i^{*})=0\right\}=1.
\end{align*}
这样，定理的证明已经完成。
\end{proof}

\end{document}
